% before: hs-hermitian-lattice
The auxiliary lattice of the argument is a module over the ring of integers of a CM field, with
a Hermitian metric at each infinite place coming from Petersson inner products. Its size is
measured by slopes, and the maps between such lattices by heights. Both depend only on the norms
at the infinite places, one embedding for each place.

% before: hs-slopes
The slope is a normalized logarithm of the inverse covolume. Its extreme values over
sublattices and quotients measure how far a lattice is from being semistable, and the
Harder--Narasimhan slopes record the whole distribution. The frame inequality below controls
the mean slope of one lattice through the distribution of the slopes of others.

% before: hs-normalization
Several normalizations of slopes are in use, and the comparison with Bost's inequality depends
on which one is meant. The one chosen here gives the ring of integers slope zero and is twice the
slope normalized by the degree of the field.

% before: hs-hn-polygon
The proof of the frame inequality works one threshold at a time: for each $\theta$ it compares
the parts of the lattices whose slopes are at least $\theta$. It needs the maximality of these
parts and the count of the slopes above $\theta$, which are standard properties of the
Harder--Narasimhan filtration.

% before: hs-bost-slope-inequality
A map of height $h$ between Hermitian lattices lowers minimal slopes by at most $h$. This is the
only arithmetic input of the frame inequality, which applies it to one map at a time.

% before: hs-almost-tight-frame
The many maps of the frame inequality are tested at a single infinite place, where they become
linear maps between complex inner product spaces. Two numbers measure how close such a family
is to a tight frame: one for the individual maps and one for the sum of their images.

% before: hs-tight-frame-meaning
The second condition of the definition is a reverse Cauchy--Schwarz inequality. Both conditions
hold with $e=\eta=0$ for an exact tight frame, and the two parameters reappear in the losses of
the frame inequality.

% before: hs-frame-complement
The linear algebra behind the frame inequality is a quantitative spanning statement. If each map
is restricted to a subspace and the subspaces have small total codimension, an almost tight
frame still fills almost all of the target, and the lemma bounds the part it misses.

% before: hs-near-orthogonal
In the application the frame is not given directly. Chapter~\ref{chapter-integral-structure}
proves Hilbert--Schmidt bounds for the blocks of the Gram matrix between the columns attached
to different words, and for the deviation of each diagonal block from the identity.

% before: hs-near-orthogonal-frame
Near orthogonality of many columns implies the two conditions of an almost tight frame. In the
application the lattices have smaller rank than the spaces of columns, so the lemma also lets
each map be restricted to a subspace, at a cost of $3\sqrt{\delta_e}$ in the tightness when the
total codimension is at most a fraction $\delta_e$.

% before: hs-share
The proof of the frame inequality must bound the number of slopes of the source lattices that
lie far below their mean. An upper bound for all of them and a lower bound for their mean
suffice, by the following form of Markov's inequality.

% before: hs-layer-cake
The frame inequality bounds, for each depth $\tau$, the number of slopes of the target lattice
below $s-\tau$. Integrating these counts over $\tau$ turns them into a lower bound for its mean
slope.

% before: hs-frame-situation
The frame inequality and its variants share a long list of data: a target lattice, many source
lattices with maps into it, bounds for their slopes and heights, and linear maps at one infinite
place. The number $s$ is a lower bound for the weighted mean of the source slopes, $u$ bounds how
far the largest of them lie above $s$, and $H$ bounds the depth of the slopes of the target below
$s$.

% before: hs-frame-slope-inequality-sharp
With these data the slope of the target lattice is bounded below by the mean slope of the
sources, up to losses controlled by the frame. The proof applies Bost's inequality to the
Harder--Narasimhan steps of each source, counts dimensions at the place $v$, and integrates over
the depth.

% before: hs-frame-slope-inequality
The application subtracts the logarithm of the number of maps as well. This weaker bound follows
at once, and it is the form recorded in the frame bound for the glued lattice.

% before: hs-comparison-with-bost
Bost's inequality and the frame inequality both bound slopes of a target lattice below by slopes
of sources and heights of maps. The comparison separates the part of the frame inequality that
comes from Bost's inequality from the part that comes from the frame, and lists the losses.

% before: hs-frame-from-near-orthogonality
Combining the frame inequality with Lemma~\ref{hs-near-orthogonal-frame} expresses its frame
hypotheses through near orthogonality, the form in which
Chapter~\ref{chapter-integral-structure} proves them. The tightness parameter becomes
$5\sigma+3\sqrt{\delta_e}$ and the isometry parameter becomes $2e_0$.

% before: hs-frame-realisation
For a definite-side datum the target is the dual $CC$ of the glued lattice $RR$, the sources
are the lattices $N_u$ attached to the words, and the maps are those of the division datum. The
datum is frame-realised when the frame data of Corollary~\ref{hs-frame-from-near-orthogonality}
exist for these maps at some infinite place, with the given parameters.

% before: hs-frame-bound
The comparison inequality of Chapter~\ref{chapter-inert-prime-balance} and the assembly in
Chapter~\ref{chapter-assembly} use the conclusion of the frame inequality for the glued lattice
in a fixed form, with the losses expressed through two parameters. The frame bound is that
conclusion, separated from the hypotheses that imply it.

% before: hs-frame-at-glued-lattice
The frame bound holds whenever the datum is frame-realised and the dimensions of the spaces
$M_u$ satisfy the provisos chosen with the scale. The provisos give at least $D^*/\sigma$
columns in total, which the near-orthogonality lemma needs.
